\chapter{Premodern Korea: From the Three Kingdoms to Joseon}

\section{Foundations}

\subsection{Myth: Dangun}

Korean national mythology begins with \textbf{Dangun} (\ko{단군}), son of the heavenly
prince Hwanung and a bear who became a woman after enduring a hundred days in a cave eating
only mugwort and garlic (the tiger gave up). Dangun founded Gojoseon, traditionally in
2333 BCE. The story is first recorded in the 13th-century \emph{Samguk yusa}, written
during the Mongol invasions, which is itself significant: national origin myths often gain
importance in times of foreign threat. October 3 is still a public holiday in South Korea
(National Foundation Day, \ko{개천절}).

\subsection{Chinese commanderies}

In 108 BCE the Han emperor Wu conquered Gojoseon and set up four commanderies; the
commandery of Lelang (Korean: Nangnang) near present-day Pyongyang lasted until 313 CE.
Lelang was a channel for Chinese writing, bronze and iron technology and luxury goods.

\section{The Three Kingdoms (c.\ 1st c.\ BCE -- 668 CE)}

\begin{center}
\begin{tikzpicture}[x=1cm,y=1cm]
  % rough peninsula outline
  \draw[thick, fill=gray!5] (-1.4,4.6) -- (2.6,5.0) -- (2.4,3.6) -- (1.6,2.6) -- (1.9,1.0)
      -- (1.6,-0.6) -- (0.4,-0.9) -- (-0.4,-0.4) -- (-0.6,0.8) -- (-0.2,1.8) -- (-0.9,2.8)
      -- (-1.4,3.4) -- cycle;
  % Goguryeo region (north, extends into Manchuria)
  \fill[red!20, opacity=0.8] (-3.2,4.0) -- (-1.4,3.4) -- (-0.9,2.8) -- (-0.2,1.8) -- (1.6,2.6)
      -- (2.4,3.6) -- (2.6,5.0) -- (1.8,6.2) -- (-2.4,6.2) -- cycle;
  \node[font=\small\bfseries, align=center] at (0.4,4.4) {Goguryeo\\\ko{고구려}};
  % Baekje (southwest)
  \fill[blue!20] (-0.2,1.8) -- (-0.6,0.8) -- (-0.4,-0.4) -- (0.4,-0.9) -- (0.6,1.9) -- cycle;
  \node[font=\small\bfseries, align=center] (bj) at (-2.0,0.9) {Baekje\\\ko{백제}};
  \draw[->] (bj.east) -- (0.0,0.6);
  % Silla (southeast)
  \fill[green!25] (0.6,1.9) -- (0.4,-0.9) -- (1.6,-0.6) -- (1.9,1.0) -- (1.6,2.6) -- cycle;
  \node[font=\small\bfseries, align=center] (sl) at (3.4,1.4) {Silla\\\ko{신라}};
  \draw[->] (sl.west) -- (1.3,1.0);
  % Gaya
  \fill[yellow!40] (0.4,-0.9) -- (0.95,-0.75) -- (0.75,0.0) -- (0.45,-0.2) -- cycle;
  \node[font=\scriptsize, anchor=east] (gy) at (-0.8,-1.3) {Gaya confederacy};
  \draw[->] (gy.east) -- (0.6,-0.55);
  % neighbours
  \node[font=\small, align=center] at (-3.4,2.6) {Yellow Sea\\(to China)};
  \node[font=\small, align=center] at (3.8,0.0) {to Japan\\(Yamato)};
  \draw[->, thick, blue!60!black] (0.0,-1.1) to[bend right=30] node[below, font=\scriptsize] {Baekje ally} (3.6,-0.6);
  \node[font=\scriptsize, gray] at (-0.3,5.8) {Manchuria};
  \node[font=\scriptsize, anchor=west] at (2.9,5.6) {c.\ 5th--6th century};
\end{tikzpicture}
\end{center}

\begin{center}
\begin{tabular}{lp{3.6cm}p{3.6cm}p{3.6cm}}
\toprule
 & Goguryeo & Baekje & Silla \\
\midrule
Location & north, into Manchuria & southwest & southeast (capital Gyeongju) \\
Character & military, horse-riding; resisted Sui and Tang invasions & cultural transmitter to Japan; maritime & last to adopt Buddhism and Chinese institutions; aristocratic ``bone-rank'' system \\
Buddhism & 372 & 384 & 527/528 \\
End & destroyed by Silla--Tang alliance, 668 & destroyed by Silla--Tang alliance, 660 & unified the peninsula south of the Taedong river \\
\bottomrule
\end{tabular}
\end{center}

\begin{example}[Goguryeo vs.\ the Sui]
The Sui emperor Yang invaded Goguryeo repeatedly (612--614) with enormous armies. The
Goguryeo general Eulji Mundeok destroyed a Sui force at the Salsu river in 612. The cost of
these failed campaigns helped bring down the Sui dynasty: an example of the peninsula
shaping Chinese history and not only the reverse.
\end{example}

\begin{aside}
Whether Goguryeo was a ``Korean'' or a ``Chinese regional'' state became a heated
diplomatic dispute in the 2000s, when a PRC academic project (the ``Northeast Project'')
classified it as part of Chinese history. History is never just about the past in this
region.
\end{aside}

\section{Unified Silla (668--935)}

\begin{itemize}
  \item Silla allied with the Tang to defeat Baekje and Goguryeo, then drove Tang forces
        out of the peninsula (676).
  \item North of Silla, Goguryeo refugees and Malgal people founded \textbf{Balhae}
        (698--926); Korean historians now speak of a ``North--South States period''.
  \item Silla's capital Gyeongju was modelled on Chang'an. Monuments: Bulguksa temple and
        the Seokguram grotto Buddha (8th c.).
  \item The \textbf{bone-rank system} (\ko{골품제}): hereditary aristocratic ranks (``sacred
        bone'', ``true bone'', head ranks 6 to 1) limited office to birth, frustrating
        talented men of lower rank, some of whom (like Choe Chiwon) studied in Tang China.
  \item The \emph{hwarang} (\ko{화랑}, ``flower youths''): aristocratic youth bands trained
        in martial arts, poetry and Buddhism; later romanticized as a Korean warrior ideal.
\end{itemize}

\section{Goryeo (918--1392)}

Founded by Wang Geon, who reunified the peninsula in 936. The name Goryeo (\ko{고려})
claims descent from Goguryeo and is the origin of the English word ``Korea''.

\begin{itemize}
  \item \textbf{Buddhism} was the state religion: the \emph{Tripitaka Koreana}; famous
        celadon ceramics (jade-green glaze with inlay).
  \item \textbf{Exams} from 958, but aristocratic families still dominated.
  \item \textbf{Military rule} (1170--1270): a military coup put generals (the Choe family)
        in power, a parallel to Japan's shoguns.
  \item \textbf{Mongol domination} (1259--1356): after decades of resistance, Goryeo became
        a ``son-in-law'' state of the Yuan; Goryeo kings married Mongol princesses, and Korean
        ships and soldiers were used in the Mongol invasions of Japan (1274, 1281).
  \item \textbf{Printing}: \emph{Jikji} (1377), oldest extant book printed with metal
        movable type.
\end{itemize}

\section{Joseon (1392--1910)}

\subsection{Founding}

General \textbf{Yi Seong-gye}, sent to attack the new Ming dynasty, turned his army around
(1388), seized power and founded Joseon in 1392 with its capital at Hanyang (Seoul). He was
backed by Neo-Confucian reformers who wanted to break the power of Buddhist monasteries and
old aristocrats.

\subsection{King Sejong and hangul}

\textbf{Sejong the Great} (r.\ 1418--1450) is the most revered Korean monarch (he is on the
10\,000 won note). His reign saw the Hall of Worthies (\ko{집현전}), advances in astronomy
(rain gauges, water clocks), agriculture manuals and above all the creation of hangul.

\begin{important}[Hangul: key facts]
\begin{itemize}
  \item Created 1443, promulgated 1446 in \emph{Hunmin jeongeum} (``Correct Sounds for the
        Instruction of the People'').
  \item Rationale (Sejong's preface): the sounds of Korean differ from Chinese, so the
        common people cannot express themselves with characters; the king created 28
        letters so that everyone could learn them easily.
  \item Opposition: officials such as Choe Malli argued that abandoning characters would
        make Korea ``barbarian'' in Chinese eyes.
  \item Status: called \emph{eonmun} (``vulgar script''), used mainly by women, commoners
        and in popular literature until the late 19th century. The name \emph{hangeul}
        (``great script'') dates from about 1912 (Ju Si-gyeong). Hangul Day: October 9.
\end{itemize}
\end{important}

\subsection{How hangul works}

Consonant letters are based on the shape of the speech organs; vowel letters are built from
three cosmological symbols: a dot (heaven, \ko{ㆍ}), a horizontal line (earth, \ko{ㅡ}) and a
vertical line (human, \ko{ㅣ}). Letters are assembled into syllable blocks.

\begin{center}
\begin{tikzpicture}[cell/.style={draw, thick, minimum size=1.1cm, font=\Large},
                    lab/.style={font=\scriptsize, align=center}]
  % 한 = ㅎ + ㅏ + ㄴ
  \node[cell, fill=blue!10] (h) at (0,1.1) {\ko{ㅎ}};
  \node[cell, fill=red!10] (a) at (1.1,1.1) {\ko{ㅏ}};
  \node[cell, fill=green!10, minimum width=2.2cm] (n) at (0.55,0) {\ko{ㄴ}};
  \node[lab, left=0.15cm of h] {initial\\h};
  \node[lab, right=0.15cm of a] {vowel\\a};
  \node[lab, below=0.1cm of n] {final n};
  \node[font=\huge] at (3.6,0.55) {$\Rightarrow$};
  \node[draw, thick, minimum size=2.2cm, font=\fontsize{40}{44}\selectfont] at (6.0,0.55) {\ko{한}};
  \node[lab] at (6.0,-0.85) {han};
  % 글 = ㄱ + ㅡ + ㄹ (vowel under)
  \node[cell, fill=blue!10, minimum width=2.2cm] (g) at (9.5,1.6) {\ko{ㄱ}};
  \node[cell, fill=red!10, minimum width=2.2cm, minimum height=0.6cm] (eu) at (9.5,0.75) {\ko{ㅡ}};
  \node[cell, fill=green!10, minimum width=2.2cm] (l) at (9.5,-0.1) {\ko{ㄹ}};
  \node[lab, right=0.15cm of g] {g};
  \node[lab, right=0.15cm of eu] {eu};
  \node[lab, right=0.15cm of l] {l};
  \node[lab] at (9.5,-1.0) {\ko{글} geul: vertical vowel stacks};
\end{tikzpicture}
\end{center}

\begin{center}
\begin{tabular}{llll}
\toprule
Basic letter & Sound & Shape represents & Derived (added stroke) \\
\midrule
\ko{ㄱ} & g/k & tongue root touching the soft palate & \ko{ㅋ} k \\
\ko{ㄴ} & n & tongue tip touching the gum ridge & \ko{ㄷ} d/t, \ko{ㅌ} t \\
\ko{ㅁ} & m & the closed mouth & \ko{ㅂ} b/p, \ko{ㅍ} p \\
\ko{ㅅ} & s & the teeth & \ko{ㅈ} j, \ko{ㅊ} ch \\
\ko{ㅇ} & (silent)/ng & the throat & \ko{ㅎ} h \\
\bottomrule
\end{tabular}
\end{center}

\begin{remark}
Linguists often call hangul a \emph{featural} alphabet: related sounds have related shapes,
because extra strokes mark aspiration or a stronger articulation. This makes it unusual
among the world's writing systems.
\end{remark}

\subsection{Neo-Confucian society}

\begin{itemize}
  \item \textbf{Yangban} (\ko{양반}, ``two orders'': civil and military officials): a
        hereditary elite that monopolized the exams and office.
  \item Below them: \emph{jungin} (technical specialists), commoners (\emph{sangmin}),
        and the lowborn (\emph{cheonmin}), including a large hereditary slave population
        (perhaps 30\% of the population in the 16th c.).
  \item Family transformation (17th c.): Goryeo-era practices (equal inheritance among sons
        and daughters, uxorilocal marriage, daughters performing rites) were replaced by
        strict patrilineal lineages, primogeniture, and ancestor rites led by the eldest
        son. Widow remarriage was stigmatized.
  \item Factional struggle among officials (``literati purges'') and powerful private
        academies (\ko{서원}, seowon).
\end{itemize}

\subsection{Invasions}

\begin{itemize}
  \item \textbf{Imjin War} (1592--1598): Toyotomi Hideyoshi invaded Korea as a step toward
        conquering Ming China. Admiral \textbf{Yi Sun-sin} and his ``turtle ships'' cut
        Japanese supply lines; Ming armies intervened; righteous armies (\ko{의병}) of
        volunteers fought on land. Devastation was enormous; Korean potters were
        carried off to Japan.
  \item \textbf{Manchu invasions} (1627, 1636): Joseon was forced to submit to the Qing.
\end{itemize}

\subsection{Late Joseon}

\begin{itemize}
  \item \textbf{Silhak} (\ko{실학}, ``practical learning''): 18th-century scholars such as
        Jeong Yak-yong criticized empty orthodoxy and studied agriculture, technology and
        administration.
  \item Kings Yeongjo and Jeongjo (18th c.) attempted to curb factionalism.
  \item 19th c.: in-law families dominated the throne; peasant revolts; the
        \emph{Donghak} (``Eastern Learning'') religious movement. These set the stage for
        Week 9.
\end{itemize}

\begin{aside}
Korean historical dramas (\emph{sageuk}) such as \emph{Dae Jang Geum} (2003) are set in
Joseon and were early hits of the Korean Wave. They are a major way that people outside
Korea (and inside it) imagine this period.
\end{aside}

\section*{Summary}
\begin{itemize}
  \item Three Kingdoms: Goguryeo (north, military), Baekje (southwest, link to Japan),
        Silla (southeast; unified the peninsula with Tang help in 668/676).
  \item Goryeo: Buddhist state, celadon, Tripitaka Koreana, Mongol domination.
  \item Joseon: Neo-Confucian state led by yangban; Sejong's hangul (1443/1446) is a
        featural alphabet arranged in syllable blocks.
  \item Joseon society became strictly patrilineal in the 17th century.
  \item Imjin War (1592--98) and Manchu invasions shaped Korean relations with Japan and
        China for centuries.
\end{itemize}
